\chapter{Metodologia}

Este capítulo descreve os procedimentos adotados na condução do trabalho. Inicialmente, apresenta-se a caracterização da pesquisa, incluindo seus objetivos e procedimentos metodológicos (pesquisa bibliográfica e aplicação de formulário). Em seguida, detalha-se o processo de engenharia de software adotado para a construção da aplicação e, por fim, o método de avaliação e validação do sistema junto aos usuários \cite{wazlawick2019}.

\section{Abordagem Metodológica}

Neste trabalho, adotou-se a abordagem de Apresentação de Produto, proposta por \citeonline{wazlawick2019}. Essa metodologia é utilizada na área da Computação por organizar a pesquisa em etapas práticas: inicia-se na identificação de um problema real e conclui-se na entrega e validação de um software funcional voltado a mitigá-lo.

Quanto à sua natureza, a pesquisa classifica-se como aplicada. Seu objetivo é conceber uma solução tecnológica direcionada a uma demanda prática: organizar e centralizar o controle de testes de software nos projetos do Instituto Federal do Espírito Santo (IFES) Campus Alegre.

Em relação aos objetivos, a pesquisa possui caráter exploratório e descritivo. A etapa exploratória ocorreu por meio do estudo da literatura sobre Garantia da Qualidade (QA) e a arquitetura MVC \cite{vasconcelos2006, delamaro2016}. Já a etapa descritiva envolveu o mapeamento dos requisitos funcionais, a modelagem do banco de dados e o desenho da arquitetura do sistema \cite{wazlawick2019}.

Por fim, a abordagem de análise combina aspectos qualitativos e quantitativos em um estudo de caso. A etapa qualitativa buscou compreender a rotina e as dificuldades das equipes de desenvolvimento do campus, permitindo que o Test-Manager fosse projetado para a hierarquia e as práticas pedagógicas do IFES, enquanto a validação empírica apoia-se na mensuração estruturada da percepção de aceitação da ferramenta pelos usuários.

\section{Processo de Desenvolvimento de Software}

Para a construção do sistema Test-Manager, adotou-se o modelo de processo de desenvolvimento \textbf{iterativo e incremental} \cite{pressman2016, sommerville2019}. Diferentemente de abordagens estritamente sequenciais (como o modelo em cascata), o desenvolvimento iterativo e incremental divide a construção do produto em ciclos curtos e evolutivos (incrementos), permitindo que as funcionalidades sejam projetadas, codificadas, testadas e refinadas continuamente a partir do retorno obtido junto ao professor orientador \cite{pressman2016, sommerville2019}.

O ciclo de engenharia do software foi estruturado em quatro incrementos principais:
\begin{enumerate}
    \item \textbf{Concepção e Infraestrutura Base:} Levantamento e elicitação de requisitos a partir da análise das planilhas de teste utilizadas no campus, configuração do ambiente conteinerizado em Docker (PHP/Laravel e MySQL) e implementação do controle de autenticação e autorização baseado em perfis (Coordenador, Tester e Desenvolvedor);
    \item \textbf{Gestão de Projetos e Casos de Teste:} Modelagem relacional e implementação dos módulos de gerenciamento de projetos, alocação de membros da equipe, cadastro manual de casos de teste e rotina de importação em lote via arquivos \texttt{.csv};
    \item \textbf{Versionamento, Execução e Rastreamento de Defeitos:} Desenvolvimento do motor de versionamento automático (\texttt{v1.0}, \texttt{v2.0}), fluxo de avaliação de testes pelos Testers, obrigatoriedade de abertura de defeitos (com anexação de evidências visuais) em caso de reprovação e ciclo de correção pelos Desenvolvedores;
    \item \textbf{Governança, Métricas e Auditoria:} Construção do painel de indicadores (\textit{dashboard}) em tempo real, geração de relatórios consolidados em PDF e implementação do log de auditoria (\texttt{Historico}) para rastreabilidade das ações.
\end{enumerate}

Em cada iteração, foram aplicadas as atividades de análise de requisitos, projeto arquitetural em UML, codificação sob o padrão MVC e testes funcionais exploratórios, buscando assegurar a estabilidade de cada módulo antes do avanço para o incremento subsequente \cite{sommerville2019}.

\section{Instrumento de Coleta de Dados e Validação}

Para avaliar a plataforma desenvolvida frente ao problema de pesquisa proposto, estruturou-se uma validação empírica fundamentada no \textbf{Modelo de Aceitação de Tecnologia} (\textit{Technology Acceptance Model} --- TAM), proposto por \citeonline{davis1989}. Concepção teórica consolidada para investigar a adoção de sistemas de informação, o modelo TAM analisa a aceitação tecnológica a partir de duas dimensões centrais \cite{davis1989, venkatesh2000}:
\begin{itemize}
    \item \textbf{Utilidade Percebida (\textit{Perceived Usefulness}):} O grau em que o usuário acredita que a utilização do sistema Test-Manager contribui para melhorar o seu desempenho, a organização e a produtividade nas rotinas de teste e gestão de defeitos \cite{davis1989};
    \item \textbf{Facilidade de Uso Percebida (\textit{Perceived Ease of Use}):} O grau em que o usuário considera que a interação com a interface da plataforma é clara, compreensível e livre de esforço excessivo \cite{davis1989}.
\end{itemize}

O instrumento de coleta de dados consiste em um questionário eletrônico estruturado no Google Forms, destinado a estudantes e coordenadores que atuam nos papéis de Coordenador, Tester e Desenvolvedor nos projetos de software do IFES Campus Alegre. As questões afirmativas do instrumento foram parametrizadas utilizando a \textbf{Escala Likert de 5 pontos} \cite{likert1932}, na qual o participante expressa seu nível de concordância segundo a gradação: (1) Discordo Totalmente; (2) Discordo Parcialmente; (3) Neutro (Nem concordo, nem discordo); (4) Concordo Parcialmente; e (5) Concordo Totalmente.

Adicionalmente, o formulário contempla questões abertas para a coleta de percepções qualitativas, sugestões de melhoria e identificação de eventuais dificuldades operacionais, permitindo relacionar os índices de aceitação com a experiência prática dos usuários avaliados.